% Generated draft: requires mathematical and source review.
\documentclass[11pt]{article}
\usepackage[margin=25mm]{geometry}
\usepackage{fontspec}
\setmainfont{Latin Modern Roman}
\usepackage{amsmath,amssymb,mathtools,unicode-math}
\setmathfont{Latin Modern Math}
\usepackage{enumitem,needspace,xurl,hyperref}
\setlength{\parindent}{0pt}
\setlength{\parskip}{5pt}
\setlength{\emergencystretch}{3em}
\newcommand{\N}{\mathbb{N}}
\newcommand{\Z}{\mathbb{Z}}
\newcommand{\Q}{\mathbb{Q}}
\newcommand{\R}{\mathbb{R}}
\newcommand{\C}{\mathbb{C}}
\newcommand{\F}{\mathbb{F}}
\newcommand{\A}{\mathbb{A}}
\newcommand{\PP}{\mathbb{P}}
\newcommand{\E}{\mathbb{E}}
\newcommand{\eps}{\varepsilon}
\newcommand{\dd}{\,\mathrm{d}}
\begin{document}
\section*{UnsolvedMath ID: 6000003}
\textbf{Source problem code:} catalogue entry
\subsection{Definitions and mathematical statement}
Let \((M,g)\) be a smooth Riemannian manifold and \(D\) its Levi--Civita connection. For a smooth totally symmetric covariant tensor \(T\in\Gamma(\operatorname{Sym}^3T^*M)\), define \(T^\sharp\) by
\[
g(T^\sharp(X,Y),Z)=T(X,Y,Z),
\]
and choose the convention
\[
\nabla_XY=D_XY-\tfrac12T^\sharp(X,Y),\qquad
\nabla_X^*Y=D_XY+\tfrac12T^\sharp(X,Y).
\]
Both connections are torsion-free and dual with respect to \(g\):
\[
Xg(Y,Z)=g(\nabla_XY,Z)+g(Y,\nabla_X^*Z).
\]
The structure is dually flat when both curvature tensors vanish, \(R^\nabla=R^{\nabla^*}=0\). Reversing the sign convention on \(T\) interchanges the two connections and does not change this existence question.

Define
\[
\mathcal D(M,g)=\{T\in\Gamma(\operatorname{Sym}^3T^*M):
 R^{\nabla(T)}=R^{\nabla^*(T)}=0\}.
\]
Amari's first question is whether \(\mathcal D(M,g)\ne\varnothing\) for every given Riemannian manifold. If there are multiple choices, the second asks to characterize the resulting dually flat structures. The source does not prescribe whether classification is of the tensors themselves, of marked connection pairs, or modulo structure-preserving isomorphisms.

Locally a dually flat structure has affine coordinates for \(\nabla\) in which \(g_{ij}=\partial_i\partial_j\phi\) for a strictly convex potential \(\phi\); this is the Hessian description. Global affine coordinates, global potentials, and completeness are stronger conditions and are not added to curvature flatness. The supplied research summary reports that not every Riemannian metric admits the construction and that local Hessian theory is substantially understood, while global and uniqueness issues are context-dependent. It provides no precise obstruction theorem or complete classification; its vague curvature-type comment is not converted into a necessary-and-sufficient curvature criterion.
\subsection{Short English statement}
Can every Riemannian metric support a pair of torsion-free flat dual connections obtained from a symmetric cubic tensor? When such tensors exist, classify the resulting structures with an explicitly chosen equivalence convention.
\subsection{Sources}
\begin{itemize}
\item Local UnsolvedMath registry, permanent ID \(6000003\), exact source code \texttt{catalogue entry}, ``Dually Flat Structures on Riemannian Manifolds.'' Proposer Shun-ichi Amari, 1998.
\item Furuhata, Matsuzoe and Urakawa, \emph{Open Problems in Affine Differential Geometry and Related Topics} (1998), item 1(c), recorded scanned PDF pages 1--2; \url{https://www.jstage.jst.go.jp/article/iis/4/2/4_2_125/_pdf/-char/en}. Recorded access 2026-07-29; manual OCR from page images.
\item Supplied PARTIAL-PROGRESS research summary invokes local Hessian characterization and curvature obstructions, but supplies no precise theorem, counterexample, or additional publication reference.
\end{itemize}
\end{document}
